% =========================================================================
% example.tex — beamerthemeDMU 主题骨架示例（内容均为占位文本）
% 编译：latexmk -xelatex example.tex   （或 xelatex 跑两遍出目录/横幅定位）
% =========================================================================
\documentclass[aspectratio=169,11pt]{beamer}
\usepackage{booktabs}
\usetheme{DMU}

% 底部进度条条目名（英文逗号分隔，与 \section 顺序一致）
\dmusetprogress{选题科学依据,已有研究不足,主要研究内容与创新点,所研究问题的难点}
% 转场页副述（英文逗号分隔，与 \section 顺序一致；不配置则无副述行）
\dmusetsectiondesc{选题依据副述占位,研究不足副述占位,内容与创新点副述占位,难点副述占位}

\title{论文题目（在此填写）}
\subtitle{硕士学位论文开题报告}
\author{报告人：姓名 \and 导师：××× 教授}
\institute{大连海事大学 · 交通运输工程学院}
\date{2026 年 9 月 28 日 · 中国大连}

% 每节自动插入过渡页（不需要时注释掉本行）
\AtBeginSection[]{\frame{\sectionpage}}

\begin{document}

% 封面：coverframe = 保留横幅、隐藏进度条与页码
\begin{coverframe}
  \titlepage
\end{coverframe}

\begin{frame}{目录}
  \begin{dmuoutline}
    \dmuoutlineitem{选题科学依据}{一句话副述占位：研究背景与需求。}
    \dmuoutlineitem{已有研究不足}{一句话副述占位：文献缺口。}
    \dmuoutlineitem{主要研究内容与创新点}{一句话副述占位：模型与方法。}
    \dmuoutlineitem{所研究问题的难点}{一句话副述占位：关键挑战。}
  \end{dmuoutline}
\end{frame}

% -------------------------------------------------------------------------
\section{选题科学依据}
% -------------------------------------------------------------------------

\begin{frame}{此页标题（占位）}
  \begin{itemize}
    \item 一级要点占位文本
      \begin{itemize}
        \item 二级要点占位文本
          \begin{itemize}
            \item 三级要点占位文本
          \end{itemize}
        \item 二级要点占位文本
      \end{itemize}
    \item 一级要点占位文本
  \end{itemize}
  \vskip2mm
  \begin{block}{色块标题（占位）}
    正文占位文本，替换为实际内容。
  \end{block}
\end{frame}

% -------------------------------------------------------------------------
\section{已有研究不足}
% -------------------------------------------------------------------------

\begin{frame}{此页标题（占位）}
  \begin{columns}[T]
    \begin{column}{0.48\textwidth}
      \begin{exampleblock}{左侧色块（占位）}
        \begin{itemize}
          \item 要点占位文本
          \item 要点占位文本
        \end{itemize}
      \end{exampleblock}
    \end{column}
    \begin{column}{0.48\textwidth}
      \begin{alertblock}{右侧色块（占位）}
        \begin{itemize}
          \item 要点占位文本
          \item 要点占位文本
        \end{itemize}
      \end{alertblock}
    \end{column}
  \end{columns}
\end{frame}

% -------------------------------------------------------------------------
\section{主要研究内容与创新点}
% -------------------------------------------------------------------------

\begin{frame}{此页标题（占位）}
  \begin{block}{表格样式（占位）}
    \centering
    \begin{tabular}{lll}
      \toprule
      列一 & 列二 & 列三 \\
      \midrule
      条目占位 & 条目占位 & 条目占位 \\
      条目占位 & 条目占位 & 条目占位 \\
      \bottomrule
    \end{tabular}
  \end{block}
  \vskip2mm
  \begin{center}
    \begin{tikzpicture}
      \draw[dmuprimary!40, dashed, fill=dmuprimary!5, rounded corners]
        (0,0) rectangle (5.2,2.4);
      \node at (2.6,1.2) {\small 插图占位（\textbackslash includegraphics）};
    \end{tikzpicture}
  \end{center}
\end{frame}

% -------------------------------------------------------------------------
\section{所研究问题的难点}
% -------------------------------------------------------------------------

\begin{frame}{此页标题（占位）}
  \begin{enumerate}
    \item 难点占位文本
    \item 难点占位文本
  \end{enumerate}
\end{frame}

\end{document}
